% GENERATED FILE. Edit canonical lessons, then run npm run generate:books.
% canonical-source-hash: fnv1a64:7be74005b443a9fe
% canonical-lessons: JA-C200-norikaeru, JA-C200-deru, JA-C200-tetsudou, JA-C200-noriba, JA-C200-teiryuujo

\chapter{Changing Trains and Waiting for Buses}
\label{ch:ja-changing-trains-and-waiting-for-buses}
\hlchaptermodality{200}

\begin{chapteropening}
\textbf{By the end of this chapter:} \emph{I can say where I change trains and wait for a bus.}

Complete the last lesson of chapter 200: I can say where I change trains and wait for a bus.
\end{chapteropening}

\section[\texorpdfstring{\ja{のりかえる} (norikaeru)}{のりかえる (norikaeru)}]{\ja{のりかえる} (norikaeru) — to change trains; to transfer}

\label{lesson:JA-C200-norikaeru}



\begin{quote}
Before the new one: say the Japanese for to come to a stop, then the Japanese for an express.
\end{quote}

\subsection*{You'll want to know: \ja{のりかえる}}
\textbf{\ja{のりかえる}} — \emph{norikaeru} — \textquotedblleft{}to change trains; to transfer\textquotedblright{}.

\begin{quote}
\textbf{\ja{ここで} \ja{ちかてつに} \ja{のりかえます。}} — \emph{koko de chikatetsu ni norikaemasu} — \textquotedblleft{}I change to the subway here.\textquotedblright{}
\end{quote}

Polite \textbf{\ja{のりかえます}}, \emph{norikaemasu}; an ichidan verb.

\begin{cousinweb}
\textbf{\ja{のりかえる}} (\emph{norikaeru}) is native Japanese (wago), built from \emph{nori}, \textquotedblleft{}riding\textquotedblright{}, and \emph{kaeru}, \textquotedblleft{}to change\textquotedblright{}, both inherited rather than borrowed (JMdict).
\end{cousinweb}

\begin{grammarlens}[title={What you've built}]
One more word for everyday life at A2.
\end{grammarlens}

\subsection*{Your turn}
\begin{itemize}
\raggedright
  \item \emph{Say it:} \emph{norikaeru}
  \item \emph{Say it:} \emph{norikaeru}, once more
  \item \emph{Say it:} \emph{koko de chikatetsu ni norikaemasu} — I change to the subway here.
  \item \emph{Recall:} say the Japanese for to come to a stop, then the Japanese for an express, then say \emph{norikaeru} again
\end{itemize}

\begin{tcolorbox}[breakable,colback=teal!4,colframe=teal!35!black,title={Before you move on}]
What does \ja{のりかえる} mean? (\textquotedblleft{}To change trains; to transfer\textquotedblright{}.) Say it once more.
\end{tcolorbox}
\section[\texorpdfstring{\ja{でる} (deru)}{でる (deru)}]{\ja{でる} (deru) — to leave; to go out}

\label{lesson:JA-C200-deru}



\begin{quote}
Before the new one: say the Japanese for an express, then the Japanese for to change trains.
\end{quote}

\subsection*{You'll want to know: \ja{でる}}
\textbf{\ja{でる}} — \emph{deru} — \textquotedblleft{}to leave; to go out\textquotedblright{}.

\begin{quote}
\textbf{\ja{とっきゅうは} \ja{もう} \ja{でました。}} — \emph{tokkyū wa mō demashita} — \textquotedblleft{}The limited express has already left.\textquotedblright{}
\end{quote}

Polite \textbf{\ja{でます}}, \emph{demasu}; an ichidan verb. Its partner \textbf{\ja{だす}}, \emph{dasu}, means \textquotedblleft{}to take out\textquotedblright{}.

\begin{cousinweb}
\textbf{\ja{でる}} (\emph{deru}) is a native Japanese word (wago), inherited rather than borrowed from Chinese; its kanji was attached later as a kun'yomi, and its older origin is not settled (JMdict, KANJIDIC2).
\end{cousinweb}

\begin{grammarlens}[title={What you've built}]
One more word for everyday life at A2.
\end{grammarlens}

\subsection*{Your turn}
\begin{itemize}
\raggedright
  \item \emph{Say it:} \emph{deru}
  \item \emph{Say it:} \emph{deru}, once more
  \item \emph{Say it:} \emph{tokkyū wa mō demashita} — The limited express has already left.
  \item \emph{Recall:} say the Japanese for an express, then the Japanese for to change trains, then say \emph{deru} again
\end{itemize}

\begin{tcolorbox}[breakable,colback=teal!4,colframe=teal!35!black,title={Before you move on}]
What does \ja{でる} mean? (\textquotedblleft{}To leave; to go out\textquotedblright{}.) Say it once more.
\end{tcolorbox}
\section[\texorpdfstring{\ja{てつどう} (tetsudō)}{てつどう (tetsudō)}]{\ja{てつどう} (tetsudō) — a railway}

\label{lesson:JA-C200-tetsudou}



\begin{quote}
Before the new one: say the Japanese for to change trains, then the Japanese for to leave.
\end{quote}

\subsection*{You'll want to know: \ja{てつどう}}
\textbf{\ja{てつどう}} — \emph{tetsudō} — \textquotedblleft{}a railway\textquotedblright{}.

\begin{quote}
\textbf{\ja{この} \ja{てつどうは} \ja{みなとまで} \ja{いきます。}} — \emph{kono tetsudō wa minato made ikimasu} — \textquotedblleft{}This railway goes as far as the harbour.\textquotedblright{}
\end{quote}

\begin{cousinweb}
\textbf{\ja{てつどう}} (\emph{tetsudō}) is Sino-Japanese (kango): two Chinese characters read with their on'yomi, readings borrowed from Chinese. They are \emph{tetsu}, \textquotedblleft{}iron\textquotedblright{} (Middle Chinese \emph{thet}) and \emph{dō}, \textquotedblleft{}way\textquotedblright{} (\emph{dawX}) (KANJIDIC2; Middle Chinese in Baxter's transcription).
\end{cousinweb}

\begin{grammarlens}[title={What you've built}]
One more word for everyday life at A2.
\end{grammarlens}

\subsection*{Your turn}
\begin{itemize}
\raggedright
  \item \emph{Say it:} \emph{tetsudō}
  \item \emph{Say it:} \emph{tetsudō}, once more
  \item \emph{Say it:} \emph{kono tetsudō wa minato made ikimasu} — This railway goes as far as the harbour.
  \item \emph{Recall:} say the Japanese for to change trains, then the Japanese for to leave, then say \emph{tetsudō} again
\end{itemize}

\begin{tcolorbox}[breakable,colback=teal!4,colframe=teal!35!black,title={Before you move on}]
What does \ja{てつどう} mean? (\textquotedblleft{}A railway\textquotedblright{}.) Say it once more.
\end{tcolorbox}
\section[\texorpdfstring{\ja{のりば} (noriba)}{のりば (noriba)}]{\ja{のりば} (noriba) — a boarding point; a bus or taxi stand}

\label{lesson:JA-C200-noriba}



\begin{quote}
Before the new one: say the Japanese for to leave, then the Japanese for a railway.
\end{quote}

\subsection*{You'll want to know: \ja{のりば}}
\textbf{\ja{のりば}} — \emph{noriba} — \textquotedblleft{}a boarding point; a bus or taxi stand\textquotedblright{}.

\begin{quote}
\textbf{\ja{のりばは} \ja{えきの} \ja{まえです。}} — \emph{noriba wa eki no mae desu} — \textquotedblleft{}The boarding point is in front of the station.\textquotedblright{}
\end{quote}

\begin{cousinweb}
\textbf{\ja{のりば}} (\emph{noriba}) is native Japanese (wago), built from \emph{nori}, \textquotedblleft{}riding\textquotedblright{}, and \emph{ba}, \textquotedblleft{}place\textquotedblright{}, both inherited rather than borrowed (JMdict).
\end{cousinweb}

\begin{grammarlens}[title={What you've built}]
One more word for everyday life at A2.
\end{grammarlens}

\subsection*{Your turn}
\begin{itemize}
\raggedright
  \item \emph{Say it:} \emph{noriba}
  \item \emph{Say it:} \emph{noriba}, once more
  \item \emph{Say it:} \emph{noriba wa eki no mae desu} — The boarding point is in front of the station.
  \item \emph{Recall:} say the Japanese for to leave, then the Japanese for a railway, then say \emph{noriba} again
\end{itemize}

\begin{tcolorbox}[breakable,colback=teal!4,colframe=teal!35!black,title={Before you move on}]
What does \ja{のりば} mean? (\textquotedblleft{}A boarding point; a bus or taxi stand\textquotedblright{}.) Say it once more.
\end{tcolorbox}
\section[\texorpdfstring{\ja{ていりゅうじょ} (teiryūjo)}{ていりゅうじょ (teiryūjo)}]{\ja{ていりゅうじょ} (teiryūjo) — a bus stop; a tram stop}

\label{lesson:JA-C200-teiryuujo}



\begin{quote}
Before the new one: say the Japanese for a railway, then the Japanese for a boarding point.
\end{quote}

\subsection*{You'll want to know: \ja{ていりゅうじょ}}
\textbf{\ja{ていりゅうじょ}} — \emph{teiryūjo} — \textquotedblleft{}a bus stop; a tram stop\textquotedblright{}.

\begin{quote}
\textbf{\ja{ていりゅうじょで} \ja{ともだちを} \ja{まちます。}} — \emph{teiryūjo de tomodachi o machimasu} — \textquotedblleft{}I wait for my friend at the bus stop.\textquotedblright{}
\end{quote}

\begin{cousinweb}
\textbf{\ja{ていりゅうじょ}} (\emph{teiryūjo}) is Sino-Japanese (kango): three Chinese characters read with their on'yomi, readings borrowed from Chinese. They are \emph{tei}, \textquotedblleft{}to halt\textquotedblright{} (Middle Chinese \emph{deng}), \emph{ryū}, \textquotedblleft{}to stay\textquotedblright{} (\emph{ljuw}) and \emph{jo}, \textquotedblleft{}place\textquotedblright{} (\emph{srjoX}) (KANJIDIC2; Middle Chinese in Baxter's transcription). The last is voiced from \emph{sho} to \emph{jo}.
\end{cousinweb}

\begin{grammarlens}[title={What you've built}]
One more word for everyday life at A2.
\end{grammarlens}

\subsection*{Your turn}
\begin{itemize}
\raggedright
  \item \emph{Say it:} \emph{teiryūjo}
  \item \emph{Say it:} \emph{teiryūjo}, once more
  \item \emph{Say it:} \emph{teiryūjo de tomodachi o machimasu} — I wait for my friend at the bus stop.
  \item \emph{Recall:} say the Japanese for a railway, then the Japanese for a boarding point, then say \emph{teiryūjo} again
\end{itemize}

\begin{tcolorbox}[breakable,colback=teal!4,colframe=teal!35!black,title={Before you move on}]
What does \ja{ていりゅうじょ} mean? (\textquotedblleft{}A bus stop; a tram stop\textquotedblright{}.) Say it once more.
\end{tcolorbox}
